\chapter{Εξόρυξη συχνών υποσυνόλων}
\label{chap:frequent-itemsets}

Η ανάλυση καλαθιού αγορών εντοπίζει προϊόντα που εμφανίζονται από κοινού στις
ίδιες συναλλαγές συχνότερα από όσο θα δικαιολογούσε η τύχη. Αποτελεί άμεση
απάντηση στο ερώτημα ποια προϊόντα παρουσιάζουν ιδιαίτερη συσχέτιση, όπως
αυτή αποτυπώνεται στα διαθέσιμα δεδομένα.

\section{Κατασκευή καλαθιών αγορών}

Ως καλάθι ορίζεται το σύνολο των διακριτών προϊόντων μιας παραγγελίας. Δύο
κανονικοποιήσεις εφαρμόζονται πριν την εξόρυξη.

Πρώτον, οι επαναλαμβανόμενες γραμμές του ίδιου προϊόντος εντός μιας
παραγγελίας συμπτύσσονται. Η συνεμφάνιση είναι δυαδική ιδιότητα: η αγορά δύο
τεμαχίων του ίδιου προϊόντος δεν προσθέτει πληροφορία συσχέτισης.

Δεύτερον, εξαιρούνται τα καλάθια με ένα μόνο προϊόν. Δεν μπορούν εξ ορισμού να
συνεισφέρουν σε καμία συνεμφάνιση, ενώ στα δεδομένα λιανικής είναι πολυάριθμα.
Η εξαίρεσή τους οξύνει τις τιμές υποστήριξης αντί να τις στρεβλώνει, καθώς ο
παρονομαστής παύει να περιλαμβάνει συναλλαγές που δεν μπορούσαν να
συνεισφέρουν.

\section{Πειραματική διάταξη}

Υλοποιούνται και οι δύο κλασικοί αλγόριθμοι. Ο \textenglish{Apriori}
\parencite{agrawal1994fast} παράγει υποψήφια σύνολα ανά επίπεδο,
αξιοποιώντας την ιδιότητα ότι κάθε υποσύνολο ενός συχνού συνόλου είναι επίσης
συχνό. Ο \textenglish{FP-Growth} \parencite{han2000mining} κατασκευάζει
δένδρο προθεμάτων και αποφεύγει εντελώς τη δημιουργία υποψηφίων.

Για το ίδιο κατώφλι υποστήριξης οι δύο αλγόριθμοι παράγουν ταυτόσημα συχνά
υποσύνολα. Η επιλογή μεταξύ τους αφορά επομένως τον χρόνο εκτέλεσης και όχι το
αποτέλεσμα, και η ισοδυναμία αυτή επαληθεύεται αυτόματα στη σουίτα ελέγχων.

Τα κατώφλια υποστήριξης και εμπιστοσύνης αποτελούν παραμέτρους της εκτέλεσης
και αποθηκεύονται μαζί με τα αποτελέσματα. Μια τιμή υποστήριξης είναι άνευ
νοήματος χωρίς το κατώφλι υπό το οποίο προέκυψε.

\section{Αποτελέσματα}

% TODO(PR7): Πίνακες αποτελεσμάτων και χρόνων εκτέλεσης μετά την εισαγωγή του
% πλήρους συνόλου δεδομένων.

\section{Αξιολόγηση κανόνων συσχέτισης}

Κάθε κανόνας αξιολογείται με τρία μέτρα. Η υποστήριξη εκφράζει το ποσοστό των
καλαθιών που περιέχουν και τα δύο σκέλη. Η εμπιστοσύνη εκφράζει την
πιθανότητα του επακόλουθου δεδομένου του προηγούμενου. Η άρση
(\textenglish{lift}) είναι ο λόγος της εμπιστοσύνης προς την υποστήριξη του
επακόλουθου.

Η άρση είναι το κρίσιμο μέτρο για την ερμηνεία. Υποστήριξη και εμπιστοσύνη από
μόνες τους αναδεικνύουν κανόνες των οποίων το επακόλουθο είναι απλώς δημοφιλές:
ένα προϊόν που αγοράζεται σχεδόν πάντοτε θα εμφανίζεται με υψηλή εμπιστοσύνη
ως επακόλουθο οποιουδήποτε προηγούμενου, χωρίς να υπάρχει πραγματική
συσχέτιση. Τιμή άρσης μεγαλύτερη της μονάδας δηλώνει ότι τα δύο σκέλη
συνεμφανίζονται περισσότερο από όσο προβλέπει η ανεξαρτησία, και για τον λόγο
αυτό το σύστημα φιλτράρει εξ ορισμού τους κανόνες με άρση κάτω της μονάδας.
